
%%% Local Variables:
%%% mode: latex
%%% TeX-master: "../main"
%%% End:

\begin{ack}
“世人希德门，揭若攀峰峦。君门起天中，多士如星攒。古来才杰士，所嗟遭时难。一鸣从此始，相望青云端。”

时间如光似箭。我从二十二岁离开中南大学，那个时候已经在长沙待了四年。今年二十五岁，在科大又待了三年，从河西左家珑搬到了黄土岭的铁道学院，后面又搬到了四方坪，长沙与我似乎有不解之缘。往后几年，要继续和这片土地联系在一起，真是又爱又恨。

二十五岁，距离二十岁和三十岁是一样远。在这个年纪是非常幸福的，快速适应AI时代的变化，对新鲜事物的接受能力处在顶峰，享受着时代进步带来的新鲜感和快感。这个年纪，不搅和到时代浪潮里去，不去折腾新东西，太可惜。还好我读了研，有了这样的机会，还要我读的还不错，能站在好平台上，还好我还年轻，可以胡闹瞎折腾。二十岁的我享受着无比幸福的大学时光，二十五岁的我还有一股劲，希望三十岁的我能像诗中所写，继续加油努力。

本文的致谢首先感谢我自己，也就是myself。感谢内容是上面的一些感想吧，感谢能读到这里。

其次感谢我的家人和女朋友。感谢她们对我的鼓励和帮助，使我能长期保持心理健康，保持积极向上的心态。尤其感谢我的女朋友谢丹妮，她时常扮演者家人的角色，我们共同克服了很多困难，也共同去探索了这个世界，这对我非常有帮助。

然后感谢我的老师、同门、舍友、队里以及科大的朋友们（为了保护他们的姓名权，不再列举姓名，都在心里）。没有他们，就难以构建起我三年的硕士生活。每个人的经历并不只有幸福，还有很多艰难的时刻。许多艰难的时刻是我们一同度过的，我认为艰难时刻的度过比幸福的时光更加来之不易，我们深厚的感情都是由此建立，并且永生难忘。

再感谢我的许多朋友（这里特指除以上以外的朋友）。他们或多或少出现在了我的时间线上，有一些很难再见到，有一些会努力地相见。我很期待与他们的下一次相见，畅谈埋藏着的记忆，诉说近来的开心与烦恼，见证彼此的成长和进步。

最后感谢这个世界。希望长沙每年的三月份不要再连下一个月的雨了，真的让人很阴郁；希望美伊不要再打仗了，二战以来世界真正和平的时间只有26天；希望AI的token费用快点降下来，真的很贵。

\end{ack}
